\section{Shortest Job First (SJF) CPU Scheduling}
\label{sec:shortest-job-first}

\problemstatement{We are given the burst times (the amount of CPU time
each process needs) of $n$ processes, all of which arrive and are ready to
run at time $0$. Only one process can run on the CPU at a time, and once we
commit to running a process we run it to completion (non-preemptive
scheduling). We must decide the order in which to run the processes so as
to \textbf{minimize the average waiting time}, where the waiting time of a
process is the amount of time it sits ready in the queue before it actually
starts executing.}

\begin{patternbox}
\emph{``Minimize average waiting time, single CPU, run-to-completion
(non-preemptive), all jobs ready at the same time''} is the textbook
\emph{Shortest Job First} scheduling rule. The keyword combination to watch
for is \textbf{``minimize the average (or total) waiting/completion
time''} together with \textbf{``one task at a time, no preemption.''}
Whenever every job is available from the very start, the order of execution
is the only choice that matters, and there is a clean, provable rule for the
best order: process the \textbf{shortest job first}. This is conceptually
close to a sorting greedy, but the insight worth absorbing is \emph{why}
ordering by burst time, rather than any other attribute, minimizes the
\emph{average} of a quantity that depends on cumulative sums.
\end{patternbox}

\subsection*{Understanding the problem carefully}

If we decide to run the processes in some order $p_1, p_2, \dots, p_n$ with
burst times $b_1, b_2, \dots, b_n$ (relabelled according to this order),
then process $p_1$ waits $0$ time units, process $p_2$ waits $b_1$ time
units (it has to wait for $p_1$ to finish), process $p_3$ waits
$b_1 + b_2$, and in general process $p_k$ waits
\[
  W_k = b_1 + b_2 + \cdots + b_{k-1}.
\]
The total waiting time is
\[
  \sum_{k=1}^{n} W_k = \sum_{k=1}^{n} \sum_{j=1}^{k-1} b_j
                      = \sum_{j=1}^{n} (n - j)\, b_j,
\]
since burst time $b_j$ is counted once for every process that comes after
it in the order -- exactly $n-j$ of them. This formula reveals the key
structural fact: \textbf{burst times that appear earlier in the order are
multiplied by a larger coefficient} ($n-1$ for the very first job, down to
$0$ for the last job). To minimize this weighted sum, we should pair the
\emph{largest} coefficients with the \emph{smallest} burst times -- which
means running the shortest jobs first.

\begin{ideabox}[Greedy Choice]
Sort the processes in increasing order of burst time, and run them strictly
in that order.
\end{ideabox}

\subsection*{Why shortest-job-first is optimal}

This is a direct instance of the \emph{rearrangement inequality}: given two
sequences of numbers, the sum of pairwise products is minimized when one
sequence is sorted increasing and the other is sorted decreasing (so that
large values are paired with small values), and maximized when both are
sorted in the same direction. Here, the coefficients $n-1, n-2, \dots, 0$
are fixed and strictly decreasing by position; to minimize
$\sum_j (n-j) b_j$, the burst times $b_j$ must be arranged in strictly
\emph{increasing} order alongside these strictly decreasing coefficients --
that is, the smallest burst time should receive the largest coefficient
(run first, waited-for by the most subsequent jobs), and the largest burst
time should receive the smallest coefficient (run last, waited-for by
nobody).

An equivalent and perhaps more intuitive exchange argument: suppose two
adjacent jobs in some schedule are out of order, i.e.\ a longer job runs
immediately before a shorter job. Swapping them leaves the waiting time of
every other job completely unchanged (the total time consumed by this pair
together is the same regardless of their internal order, so everything
after them starts at the same time either way), but it strictly decreases
the combined waiting time of \emph{this pair}, because the shorter job, now
running first, makes the longer job wait less than the longer job, when
running first, made the shorter job wait. Repeatedly fixing such
adjacent inversions -- exactly what sorting does -- can only decrease total
waiting time, so the fully sorted (shortest first) order is optimal.

\subsection*{Algorithm}

\begin{enumerate}
  \item Sort the burst times in increasing order.
  \item Initialize $\textit{waitingTime} \gets 0$ and
        $\textit{totalWait} \gets 0$.
  \item For $i \gets 0$ to $n-1$ (processing jobs in sorted order):
  \begin{enumerate}
    \item Add $\textit{waitingTime}$ to $\textit{totalWait}$ (this is the
          waiting time of the current job).
    \item $\textit{waitingTime} \gets \textit{waitingTime} +
          \textit{burst}[i]$ (the next job must wait for this one to
          finish too).
  \end{enumerate}
  \item Return $\textit{totalWait} / n$ as the average waiting time.
\end{enumerate}

\begin{algorithm}[H]
\caption{Shortest Job First Scheduling}
\begin{algorithmic}[1]
\Procedure{SJF}{\textit{burst}}
  \State sort \textit{burst} in increasing order
  \State $\textit{waitingTime} \gets 0$, $\textit{totalWait} \gets 0$
  \For{$i \gets 0$ to $n-1$}
    \State $\textit{totalWait} \gets \textit{totalWait} + \textit{waitingTime}$
    \State $\textit{waitingTime} \gets \textit{waitingTime} + \textit{burst}[i]$
  \EndFor
  \State \Return $\textit{totalWait} / n$
\EndProcedure
\end{algorithmic}
\end{algorithm}

\subsection*{Dry run}

\dryrun{Take burst times $[6, 1, 8, 9, 3]$.}

Sorted increasing: $[1, 3, 6, 8, 9]$.

\begin{itemize}
  \item $i=0$ (burst $1$): job waits $\textit{waitingTime}=0$.
        $\textit{totalWait}=0$. Update $\textit{waitingTime}=0+1=1$.
  \item $i=1$ (burst $3$): job waits $1$. $\textit{totalWait}=1$. Update
        $\textit{waitingTime}=1+3=4$.
  \item $i=2$ (burst $6$): job waits $4$. $\textit{totalWait}=5$. Update
        $\textit{waitingTime}=4+6=10$.
  \item $i=3$ (burst $8$): job waits $10$. $\textit{totalWait}=15$. Update
        $\textit{waitingTime}=10+8=18$.
  \item $i=4$ (burst $9$): job waits $18$. $\textit{totalWait}=33$. Update
        $\textit{waitingTime}=18+9=27$.
\end{itemize}

Total waiting time is $33$, and the average waiting time is
$33 / 5 = 6.6$.

\subsection*{C++ implementation}

\begin{lstlisting}
#include <bits/stdc++.h>
using namespace std;

double shortestJobFirst(vector<int>& burst) {
    sort(burst.begin(), burst.end());

    int n = burst.size();
    long long waitingTime = 0, totalWait = 0;

    for (int i = 0; i < n; i++) {
        totalWait += waitingTime;
        waitingTime += burst[i];
    }

    return (double)totalWait / n;
}

int main() {
    vector<int> burst = {6, 1, 8, 9, 3};
    cout << "Average waiting time: "
         << shortestJobFirst(burst) << endl;
    return 0;
}
\end{lstlisting}

\begin{complexitybox}
\textbf{Time Complexity.} Sorting the $n$ burst times costs
$O(n \log n)$. The subsequent single pass to accumulate waiting times costs
$O(n)$. The overall time complexity is
\[
  O(n \log n).
\]
\textbf{Space Complexity.} Sorting is performed in place, and we use only a
constant number of extra scalar variables, giving auxiliary space
complexity $O(1)$ (ignoring sort's recursion stack).
\end{complexitybox}

\begin{takeawaybox}
Whenever the objective is a sum of the form $\sum (\text{fixed
position-dependent coefficient}) \times (\text{value you get to order})$,
and the coefficients are monotonic in position, the rearrangement
inequality tells you exactly how to order the values for the minimum (or
maximum) possible sum: pair the largest coefficients with the smallest
values (or vice versa for maximization). Recognizing this pattern --
``minimize a sum of weighted, orderable quantities'' -- immediately tells
you to sort, without needing to reconstruct the exchange argument from
scratch every time.
\end{takeawaybox}
